\chapter*{Список сокращений и условных обозначений}             
\addcontentsline{toc}{chapter}{Список сокращений и условных обозначений}  

\textbf{ИМК} --- интерфейс мозг-компьютер.

\textbf{МЭГ} --- магнитоэнцефалография.

\textbf{ЭКоГ} --- электрокортикография.

\textbf{ЭЭГ} --- электроэнцефалография.

\textbf{BSS} (Blind Source Separation) --- слепое разделение источников.

\textbf{CCA} (Canonical Correlation Analysis) --- канонический корреляционный анализ.

\textbf{CorrCA} (Correlated Component Analysis) --- анализ коррелированных компонент.

\textbf{cSPoC} (Canonical Source Power Correlation Analysis) --- канонический анализ корреляции мощности источников.

\textbf{ECCA / Envelope-CorrCA} (Envelope Correlated Component Analysis) --- анализ коррелированных компонент огибающей.

\textbf{ERD} (Event-Related Desynchronization) --- событийно-связанная десинхронизация.

\textbf{ERP} (Event-Related Potentials) --- событийно-связанные (вызванные) потенциалы.

\textbf{ERS} (Event-Related Synchronization) --- событийно-связанная синхронизация.

\textbf{ICA} (Independent Component Analysis) --- анализ независимых компонент.

\textbf{IVA} (Independent Vector Analysis) --- анализ независимых векторов.

\textbf{MDM} (Minimum Distance to Mean) --- классификатор минимального расстояния до среднего.

\textbf{mSPoC} (Multivariate Source Power Comodulation) --- многомерная комодуляция мощности источников.

\textbf{PCA} (Principal Component Analysis) --- метод главных компонент.

\textbf{Power-CCA} --- канонический корреляционный анализ мощности.

\textbf{RAP-MUSIC} (Recursively Applied and Projected Multiple Signal Classification) --- метод рекурсивной пространственной фильтрации и классификации множественных сигналов.

\textbf{SNR} (Signal-to-Noise Ratio) --- отношение сигнал/шум.

\textbf{SPD-матрица} (Symmetric Positive Definite) --- симметричная положительно определенная матрица.

\textbf{SPoC} (Source Power Comodulation) --- комодуляция мощности источников.

\textbf{SSD} (Spatio-Spectral Decomposition) --- пространственно-спектральная декомпозиция.

\textbf{STFT} (Short-Time Fourier Transform) --- оконное преобразование Фурье.

\textbf{TRCA} (Task-Related Component Analysis) --- анализ компонент, связанных с задачей.

\textbf{t-SNE} (t-distributed Stochastic Neighbor Embedding) --- стохастическое вложение соседей с t-распределением.

\textbf{UMAP} (Uniform Manifold Approximation and Projection) --- аппроксимация и проекция равномерного многообразия.